\section{Related Work}
\label{sec:related}

\paragraph{Classical index advisors.}
Index selection as an optimization problem dates to the AutoAdmin project
\citep{chaudhuri1997autoadmin}, which searches over candidate index
configurations using a what-if optimizer. The DB2 Advisor
\citep{valentin2000db2advisor} and the DB2 Design Advisor
\citep{zilio2004db2design} extended this to multiple design structures and budget
constraints. These systems assume a fixed workload, solve an offline
optimization problem, and are re-run manually when the workload changes.

\paragraph{Learned database components.}
Machine learning now pervades the database internals. Learned index structures
\citep{kraska2018case} replace B-trees with models; learned cardinality
estimators \citep{kipf2019learnedcardinalities} and learned cost models
\citep{marcus2019learnedcost} improve plan costing; and Bao \citep{marcus2021bao}
learns to steer the query optimizer at plan time. Each of these components is
trained offline on a corpus and then frozen; none addresses the question of how a
\emph{deployed} model should keep learning as the workload evolves.

\paragraph{Reinforcement learning for physical design.}
A growing body of work casts physical design as an RL problem. In the bandit
setting, an agent chooses configuration actions (create/drop an index,
materialize a view) and receives a latency-based reward. Such approaches are
attractive because they avoid an explicit what-if oracle, but they inherit the
stationarity assumption: the agent is trained and evaluated on a fixed workload
distribution, and a distribution shift degrades performance. To the best of our
knowledge, CLDB is the first to treat index selection as a \emph{continual} RL
problem and to combine continual-learning machinery with an index-aware RL
formulation.

\paragraph{Continual learning.}
Continual learning (CL) seeks to learn a sequence of tasks without forgetting.
Regularization methods constrain updates to protect parameters important for
previous tasks; elastic weight consolidation (EWC) \citep{kirkpatrick2017ewc}
does so using the Fisher information matrix, while learning-without-forgetting
\citep{li2017lwf} distils the old model's outputs. Rehearsal methods store and
replay a subset of past data, either uniformly \citep{rolnick2019replay} or
through gradient episodic memory \citep{lopezpaz2017gem}, which projects
gradients to avoid increasing the loss on buffered examples. The last two
references also define the ACC/BWT/FWT metrics we adopt. Our contribution is not
a new CL algorithm but the \emph{combination} of these ideas with a
database-specific priority signal and a drift-triggered consolidation rule taken
from the concept-drift literature \citep{gama2014drift}.

\paragraph{Continual and non-stationary RL.}
Reinforcement learning under non-stationarity is surveyed by
\citet{khetarpal2022continual}. Deep RL relies on experience replay and target
networks \citep{mnih2015dqn,lillicrap2016ddpg}; prioritized replay
\citep{schaul2016prioritized} reweights transitions by TD error with importance
sampling; double Q-learning \citep{vanhasselt2016double} reduces overestimation.
Applying these ideas to a stream of \emph{conflicting} tasks is exactly the
continual-RL setting, and EWC was originally demonstrated on sequential Atari
games \citep{kirkpatrick2017ewc}. CLDB instantiates that setting for index
selection.

\paragraph{Gap.}
Existing index advisors are stationary; existing learned-DB components are frozen
after training; and existing continual-learning methods are demonstrated on
generic benchmarks rather than database tuning. CLDB sits at the intersection of
the three: a \emph{continual}, \emph{index-aware}, \emph{safety-constrained}
reinforcement-learning advisor, evaluated with continual-learning metrics on a
reproducible workload-drift benchmark.
